\section{Circuits on an open chain placed in a longer open chain}
\label{sec:circuitplacement_chain}

\begin{lemma}[A circuit on an open chain placed at arbitrary sites]
    \label{lem:chain_circuit_placement}
    \lean{QuantumCircuit.IsGateProduct.placeAlong_pathGraph,
        QuantumCircuit.IsGateProduct.placeAlong_pathGraph_of_succ}
    \leanok
    \uses{thm:unitary_swap_routing, thm:unitary_relabel, lem:sites_place_along_hom}
    Let $U$ be a product of $K$ gates on neighbouring sites of the open chain of
    $m$ sites, with the same finite local basis at every site, and let
    $w$ be an injection of the sites $\{0,\dots,m-1\}$ into those of the open
    chain of $N$ sites. Then $U$ acting on the sites $w(0),\dots,w(m-1)$, with
    the identity on the remaining sites, is a product of at most $2NK$ gates on
    neighbouring sites of the long chain. If moreover $w(p+1)=w(p)+1$ for every
    $p$, it is a product of at most $K$ such gates.
\end{lemma}

\begin{proof}\leanok
    Every gate of the short chain acts on two neighbouring sites $p,p+1$, so its
    placement acts on the two distinct sites $w(p),w(p+1)$ of the long chain.
    By swap routing, a unitary on two sites at distance $j$ is a product of at
    most $2j-1\leq 2N$ gates on neighbouring sites. When $w$ carries
    neighbouring sites to neighbouring sites, the placed gates are already gates
    on neighbouring sites, and relabelling the sites preserves the gate count.
\end{proof}

\begin{definition}[Placement of a register with its workspace]
    \label{def:clean_circuit_placement_sites}
    \lean{QuantumCircuit.cleanImplementationSites,
        QuantumCircuit.cleanImplementationSites_castAdd,
        QuantumCircuit.cleanImplementationSites_natAdd}
    \leanok
    Consider a register of $m$ logical sites followed by $b$ scratch sites, and
    a larger register of $n$ logical sites followed by $A$ scratch sites. Given
    injections $e$ of the $m$ logical sites into the $n$ logical sites and $t$ of
    the $b$ scratch sites into the $A$ scratch sites, the placement of the small
    register sends its logical site $j$ to $e(j)$ and its scratch site $j$ to
    $n+t(j)$.
\end{definition}

\begin{lemma}[Clean placement into a larger common workspace]
    \label{lem:clean_circuit_placement}
    \lean{QuantumCircuit.IsCleanImplementation.placeAlong,
        QuantumCircuit.isGateProduct_isCleanImplementation_placeAlong}
    \leanok
    \uses{def:unitary_clean_implementation, def:clean_circuit_placement_sites,
        lem:chain_circuit_placement}
    Let $U$ act on $m$ logical and $b$ scratch sites and cleanly implement $Z$
    on the logical sites with the scratch sites in the zero configuration:
    $U(\ket{x}\otimes\ket{0})=(Z\ket{x})\otimes\ket{0}$ for every logical
    configuration $x$. Then $U$ placed along the placement of the register
    cleanly implements $Z$ placed along $e$, with all $A$ scratch sites of the
    larger register in the zero configuration. If $U$ is a product of $K$ gates
    on neighbouring sites of the open chain of $m+b$ sites, its placement is a
    product of at most $2(n+A)K$ gates on neighbouring sites of the open chain of
    $n+A$ sites.
\end{lemma}

\begin{proof}\leanok
    On a configuration whose scratch part vanishes, the placed matrix acts by
    $U$ on the restriction to the placed sites and by the identity elsewhere,
    and the restriction of such a configuration has a vanishing scratch part. So
    the entries of both sides of the placed identity are the entries of the
    local clean identity, multiplied by the condition that the two
    configurations agree off the placed sites and vanish on the scratch sites
    outside the image of $t$. The gate count is the first part of
    Lemma~\ref{lem:chain_circuit_placement}.
\end{proof}
